%% kapa: exam=1 round=26 year=2015 period=2 form=A subject=appraisal_law subject_ko=부동산관계법규 seq=7 number=7
%% source: Q-Net 감정평가사 기출문제 (cst00309) artlSeq=5138001 file=extracted/1008/5138001/제26회_감정평가사 1차_기출/제26회_감정평가사 1차_기출/2교시_감평_A형.pdf
%% conversion: kordoc 4.24.51 -> tools/gen1.py (auto); review=proofread-sonnet; flags=none
\begin{question}{7}
국토의 계획 및 이용에 관한 법령상 계획관리지역의 건폐율 및 용적률의 최대한도를 바르게 연결한 것은? (단, 조례는 고려하지 않음)
\begin{choices}
  \item 건폐율: 20\% 이하, 용적률: 80\% 이하
  \item 건폐율: 20\% 이하, 용적률: 100\% 이하
  \item 건폐율: 40\% 이하, 용적률: 80\% 이하
  \item 건폐율: 40\% 이하, 용적률: 100\% 이하
  \item 건폐율: 50\% 이하, 용적률: 120\% 이하
\end{choices}
\end{question}
\answer{4}
